\documentclass[11pt, oneside]{article} 
\usepackage{geometry}
\geometry{letterpaper} 
\usepackage{graphicx}
	
\usepackage{amssymb}
\usepackage{amsmath}
\usepackage{parskip}
\usepackage{color}
\usepackage{hyperref}

\graphicspath{{/Users/telliott/Dropbox/Github-Math/figures/}}
% \begin{center} \includegraphics [scale=0.4] {gauss3.png} \end{center}

\title{Equilateral triangles}
\date{}

\begin{document}
\maketitle
\Large

%[my-super-duper-separator]

\subsection*{basic triangle}

As we saw previously, an equilateral triangle has all three sides the same length.  By symmetry, as there is no reason to favor one vertex, the three vertices have the same angular measure, namely, $\pi/3$ (since the total is two right angles).

\begin{center} \includegraphics [scale=0.3] {equilateral.png} \end{center}

\emph{Proof}.  We also have \hyperref[sec:Euclid19]{\textbf{Euclid $I.19$}}, which says that if one angle is larger than another in a triangle, then the corresponding side is also longer.  So, assume the angles are not the same, then that implies the sides are not the same length either.  Hence the angles are the same.  $\square$

An equilateral triangle is also isosceles (three times over), and as a result the altitude dropped from a vertex to the opposing base bisects both that base and the angle at the vertex.

\begin{center} \includegraphics [scale=0.3] {iso13.png} \end{center}

\subsection*{area}

To calculate the area of the triangle, we could write a general formula, but it is usually easier to go back to basics and derive one.

\begin{center} \includegraphics [scale=0.4] {pi_calc5.png} \end{center}

In the left panel, the result of bisection is that, for a side of length $2$, half the base is $1$, and then Pythagoras tells us that the altitude is $\sqrt{3}$.  Using the standard formula, the area of this triangle is also $\sqrt{3}$.

To re-scale the triangle, write whichever measurement is given.  For example in the middle we're given a side length of $1/2$ so that means that all lengths are multiplied by $1/4$.  On the right, we're given that the altitude is $1/2$ so then we must multiply by $1/2 \sqrt{3}$.

The logic is that to erase that $\sqrt{3}$ in the altitude, we must multiply by its inverse and also another factor that gets us where we need to be.

If you really insist on a general formula...

If the side is of length $s$, the length of half the base is $s/2$ and Pythagoras gives us the altitude's length as
\[ \sqrt{s^2 - \frac{s^2}{2^2}} = \frac{\sqrt{3}}{2} s \]

Therefore the area of the triangle is
\[ A = \frac{1}{2} \cdot \frac{\sqrt{3}}{2}s \cdot s = \frac{\sqrt{3}}{4} s^2 \]

We can check that with a general formula for the area of any triangle.  The area is one-half the product of two adjacent sides times the sine of the angle between.  Here the angle is $\pi/3$ and its sine is $\sqrt{3}/2$ so the area is $\sqrt{3}/4 s^2$, as we said.

It can be quicker to rely on the fact that the area goes like the square of the side.  For side length $2$ we had $A = \sqrt{3}$.  If we shrink the triangle to side length $1$, the area goes down by a factor of four, to $A = \sqrt{3}/4$.

\subsection*{circumscribed equilateral triangle}

Here is a fun construction based on an equilateral triangle.  

Any triangle fits into a unique circle (see above).
\begin{center} \includegraphics [scale=0.4] {one_third.png} \end{center}

Draw the radius to a vertex of the triangle, and extend it as the diameter.  We will show that all three  marked angles are equal.

\emph{Proof}.  

The altitude of an isosceles (or equilateral) triangle is also a median (it bisects the base) and an angle bisector (it bisects the top angle), so it goes through the center of the circle containing all three vertices.

This accounts for the two angles labeled $\theta$.  Since the sum of angles in a triangle is equal to two right angles, each of the vertex angles in an equilateral triangle is equal to $2/3$ of a right angle.  $\theta$ is one-half that measure, or $1/3$ of a right angle.

Any triangle with the diagonal for one side and the opposing vertex anywhere else on the circle is a right triangle.  That's our fundamental theorems about circles.

As a result,  $\phi$ measures one right angle minus $2/3$ of a right angle or $1/3$ of a right angle.  Therefore, $\phi = \theta$.

\begin{center} \includegraphics [scale=0.4] {one_third.png} \end{center}

So now we have one of the complementary angles of a right triangle, and a shared side.  The two triangles are congruent.  That accounts for the duplicated $a$ in the figure.

$\square$

Thus, the altitude of the equilateral triangle is $3/4$ of the diameter of the circle that just encloses it.  And the point where the altitudes meet in an equilateral triangle is $1/3$ of the way up from the base, since $r = 2a$.  

Another important result about circles is that if two angles with each vertex on the circle, subtend the same arc of the circle, then they are equal.  Since $\phi$ subtends the same arc as one-half of the top vertex of the triangle, this also shows that $\phi$ is equal to $\theta$.

\subsection*{Bertrand's paradox}

Grinstead and Snell's wonderful \emph{Introduction to Probability} has this problem (example 2.6).  It's called Bertrand's paradox.  We are asked to draw a chord of a unit circle randomly.

\begin{center} \includegraphics [scale=0.6] {Bertrand1.png} \end{center}

Here we might say, let's choose each of three angles $\alpha$, $\beta$ and $\theta$ randomly (uniform density) from $[0, 2 \pi]$.  

But there is no reason why the radius to $B$ cannot lie along the $x$-axis, so there are really only two choices.  

The question is posed:  what is the probability that the length of this random chord is $> \sqrt{3}$.

However, there are several different approaches to parametrize the problem, and randomizing the different parameters leads to different results.

\subsection*{equilateral triangles}

We briefly review some properties of equilateral triangles which we looked at earlier, with a slightly different take.

Drop an altitude and observe the ratio of side lengths.  It is convenient to start with a side length of 2 for the original triangle, then in the bisected copies the sides are in the ratio 1-2-$\sqrt{3}$ (the $1$ by bisection, and $\sqrt{3}$ by the Pythagorean theorem).

The angle at each vertex of the original equilateral triangle is $\pi/3$, so the new triangles have angles of $\pi/6$, both by bisection or because the altitude forms an angle of $\pi/2$ at the base, so the sum of angles theorem gives us the last angle.

In the right panel, the equilateral triangle is inscribed in a unit circle, so $OR = OP = OQ = 1$.  We claim that the line segment $OM$ has a length of $1/2$.

\begin{center} \includegraphics [scale=0.6] {Bertrand2.png} \end{center}

\emph{Proof}.

$\angle PQR$ is a right angle, by Thales' circle theorem, and $\angle MRQ$ is shared, so $\triangle PQR$ is similar to $\triangle RMQ$.  Therefore, $\angle RPQ = \angle MQR = \pi/3$.

Therefore the sides of $\triangle PQR$ are also in the ratio 1-2-$\sqrt{3}$, with $PQ/PR = 1/2$ and so $PQ = OP = OQ$.  Thus, $\triangle OPQ$ is equilateral.

$QM \perp OP$ so $MQ$ is the bisector of both $\angle PQO$ and the base $OP$.  

Therefore, $OM$ is one-half of $OP$ and has a length of $1/2$.

$\square$

\subsection*{first parametrization}

We have just shown that the altitude of the inscribed equilateral triangle in a unit circle has length $3/2$.  This means that the ratio of the inscribed triangle to the standard one is $\sqrt{3}/2$.

And that means the side length of the inscribed equilateral triangle is $\sqrt{3}$.

That explains the length chosen for the chord in this problem.  We see that if $M$ is chosen at random anywhere along $OP$, one-half of the time the chord formed will be larger than $\sqrt{3}$.

\begin{center} \includegraphics [scale=0.5] {Bertrand2.png} \end{center}

So the probability we are asked to give is just $1/2$.

\subsection*{second parametrization}

\begin{center} \includegraphics [scale=0.5] {Bertrand3.png} \end{center}

The second parametrization has the same triangle we just saw, $\triangle PQR$.  The angle at vertex $R$ is $\theta$.

$\theta$ can lie in the interval $[0, \pi/2]$, and in the event that $\theta < \pi/6$, the chord length $RQ > \sqrt{3}$.

The probability that the chord is greater than $\sqrt{3}$ in length is $1/3$, since $\pi/6$ is one-third of $\pi/2$.

\subsection*{third parametrization}

Finally, we imagine picking two coordinates $(x,y)$ at random from the interior of the circle.  We place the midpoint of the chord $M$ at $(x,y)$.

If $M$ is such that $r = \sqrt{x^2 + y^2} < 1/2$, then $M$ will be closer to the center of the circle than $1/2$ and so the chord length will be $> \sqrt{3}$.

The number of points that have this property is proportional to the relative areas of the inside small circle, and the larger circle around is.

\[ \frac{\pi (1/2)^2}{\pi} = \frac{1}{4} \]

We see that, depending on which parameter is randomized, we obtain a probability of $1/2$, $1/3$ or $1/4$.

In Jaynes' words, the problem is not well-formed.

\subsection*{problem}

Here is a problem from Brown, 1901.

\begin{center} \includegraphics [scale=0.4] {Brown1901_4.png} \end{center}

A regular triangle is what we would call an equilateral triangle.  Applying the Pythagorean theorem to one-half of such a triangle yields 1-2-$\sqrt{3}$ for the sides, where $2$ is the hypotenuse.

\begin{center} \includegraphics [scale=0.5] {Brown1901_4b.png} \end{center}

From this we deduce that the altitude is $3r$ and the side length is $2 \sqrt{3} r$, so the triangle's area is 
\[ A = \frac{1}{2} \cdot 3 \cdot 2 \sqrt{3} \cdot r^2 \]

The red part is the difference between this and the circle's area.

\[ \text{red} = (3 \sqrt{3} - \pi) r^2 \]

We're given a radius of $7$ but that's arithmetic.  We don't need that now.

\subsection*{problem}

Here is a problem from the entrance exam for Harvard, 1901.

\begin{center} \includegraphics [scale=0.4] {harvard1901_6.png} \end{center}

The first thing to notice is that the circle is tangent to the sides.  The tangent touches the circle at a single point, so the missing part of the diagram looks like this:

\begin{center} \includegraphics [scale=0.6] {harvard1901_6a.png} \end{center}

The whole angle at the top is $\pi/3$ so the half-angle is $\pi/6$.  This right triangle has sides in the ratio 1-2-$\sqrt{3}$.  The radius is given as $1$ so the red base side of the triangle is $\sqrt{3}$.  

The missing perimeter of the triangle is twice this or $2 \sqrt{3}$.  But the added perimeter from the circle is $2/3 \cdot 2 \pi r = 4/3 \pi$.

The area of the triangle is $1/2 \cdot \sqrt{3} \cdot 1 = \sqrt{3}/2$ and there are two copies so the total missing triangular area is $\sqrt{3}$.  The added circular area is $2/3 \pi$. 

The original perimeter of the triangle was $7 \cdot 3 = 21$ and the area was $\sqrt{3}/4 \cdot s^2$.

We'll just set up the two calculations:

\[ A = \frac{\sqrt{3}}{4} \cdot 7^2 - \sqrt{3} + 2/3 \cdot \pi \]
\[ P = 7 \cdot 3 - 2 \sqrt{3} + 4/3 \cdot \pi \]

It seems that the most efficient way to calculate this is to do $ \sqrt{3} + 2/3 \cdot \pi$, then double it, then do the rest.

\end{document}